% =========================================================================
% 5. DISCUSSION
% =========================================================================
\section{Discussion}
\label{sec:discussion}

\paragraph{The measurement problem extends beyond the four evaluated methods.}
The clearest practical result of \S\ref{sec:compare} is that an absolute
harmful-answer rate can reorder interventions and attribute pre-existing model
harm to the intervention. Baseline-corrected Potency removes this avoidable
confound, while Coherence Retention distinguishes selective safety failure from
broad response degradation. We therefore recommend that white-box intervention
studies report an induced-effect measure and a coherence or utility measure
alongside raw attack success, even when they do not adopt $\SFS$. The full
$(\Ppot,\Qual,\Sreason)$ vector should remain the primary report: the Pareto view
in Figure~\ref{fig:appendix-pareto} makes trade-offs explicit, whereas $\SFS$
provides a compact headline summary. These corrections do not eliminate
differences in datasets, judges, or intervention access, but they make
cross-model and cross-study comparisons more interpretable.

\paragraph{The tested families provide unequal behavioral evidence for localization.}
Across the safety-trained arms, direction-based methods produce the strongest
selective-failure signatures on average, but the ordering is model-dependent
and should not be interpreted as a matched-strength comparison. Steering is
strongest on \model{Llama-3.1-8B} and \model{OLMo-3-Think}, directional
ablation is strongest on \model{Qwen3-8B}, and large-budget neuron ablation has
a substantial effect on Llama. By contrast, head and neuron ablation are often
weak or accompanied by coherence loss. The pre-alignment controls further show
that target-discovery procedures can return ranked targets on a checkpoint with
little alignment-induced safety, while raw harmful-answer rates can make those
interventions appear effective even when baseline-corrected Potency is small.
These results grade the operational localization claims supported by the tested
selectors and budgets. A high $(\Ppot,\Qual,\Sreason)$ signature provides causal
behavioral evidence consistent with localization, but it cannot establish that
a target is the unique, minimal, or sufficient safety mechanism; a downstream
or correlated substrate could produce the same signature.

\paragraph{The visible pathway shifts are consistent with downstream bypass.}
In the explicit-trace models, harmful completions often preserve visible harm
recognition and refusal initiation, while execution and operational details
increase. This pattern is most pronounced under steering, which produces the
largest pathway shift in our grid. Under the tested conditions, the result is
consistent with steering changing how visible safety considerations influence
the final answer, rather than simply erasing their appearance in the trace. This
is a trace-level inference, not a localization of the underlying latent
computation. It neither rules out a component-level implementation nor implies
that safety is literally represented by a single direction: a residual-stream
direction may aggregate features implemented across heads or neurons, and weak
component-ablation results may reflect limitations of the selector or budget
rather than the irrelevance of those components.

\paragraph{Monitorability is a diagnostic property, not an intervention objective.}
Monitorability Retention does not make an intervention intrinsically more
desirable; it records whether intervention-induced harmful behavior becomes
less detectable from the available trace. Across the $100$ intervention cells,
only four have a positive signed monitorability gap, all of small magnitude,
and none is a high-Potency cell. The negative tail instead reflects
over-warning, particularly when coherence-degrading interventions leave a
safety-sensitive trace or response prefix but produce an unusable answer. On
the explicit-trace models, stronger steering shifts the pathway signature
toward execution while recognition and refusal remain visible, making the trace
more conflicted rather than more sanitized. This is bounded but encouraging
evidence that the tested interventions do not systematically evade a
trace-only monitor \citep{baker2025cotmonitoring,korbak2025monitorability}; it
is not a general guarantee for chain-of-thought monitoring. Only two arms expose
explicit traces, visible traces need not faithfully represent latent
computation, and a binary prediction of final harmfulness can miss
fine-grained process changes. Monitoring pathway-level behaviors may therefore
be more informative than relying on outcome prediction alone.

\paragraph{Limitations.}
\label{sec:limitations}
Several caveats bound these conclusions. First, $\Qual$ measures coherence
rather than benign capability: an intervention could preserve fluent responses
while undermining task competence, and the coherence gate lacks dedicated human
validation. Second, the headline harmfulness and trace-monitor judges attain
moderate human agreement, while the pathway judge's strongest validation is
transfer-domain; the fine-grained pathway labels should therefore remain
descriptive. The monitorability finding is additionally constrained by the
moderate reliability of the trace-only predictor and by the fact that only two
arms expose explicit traces; the early-response-prefix results for the other
arms are sensitivity analyses rather than chain-of-thought evidence. Third, the
grid covers five 7–8B dense, text-only model/configuration arms, a single
generation seed, and inference-time interventions with family-specific,
non-commensurate dose ladders. Our intervals bootstrap prompts rather than
model seeds or repeated generations. The family patterns therefore require
verification at larger scales, across random seeds, on mixture-of-experts and
multi-modal models, and for fine-tuning-based interventions. Finally, the two
pre-alignment \model{OLMo-3-Base} arms diagnose baseline confounding and target
provenance; they neither prove that the checkpoint contains no safety-relevant
computation nor provide a direct faithfulness ranking between methods.

\paragraph{Future work.}
The most direct mechanistic extension is a bidirectional test: a target whose
removal weakens refusal should also be tested for whether the inverse
intervention strengthens safety. The same framework can evaluate this setting
by reversing the orientation of $\Ppot$ while retaining $\Qual$ and $\Sreason$.
Coherence Retention should also be extended with benign-utility measurements,
and monitorability should be tested on models with compressed, latent, or
otherwise non-verbalized reasoning. Model-level robustness should not be
inferred from grid-average susceptibility until the intervention portfolio and
budgets are standardized; a useful benchmark would calibrate these choices and
estimate variation across seeds and generations. Finally, combining
direction-level interventions with component-level causal analysis could test
whether successful directions aggregate distributed features or are
implemented by identifiable circuits.
